\hypertarget{local-development-setup}{%
\section{Local Development Setup}}

This section covers setting up the complete local development
environment using Docker Compose to simulate the AWS services locally.

\hypertarget{architecture-overview-local}{%
\subsection{Architecture Overview
(Local)}}

\emph{(Mã nguồn chi tiết đã được lược bỏ để báo cáo ngắn gọn. Vui lòng
xem trong source code dự án)}

\hypertarget{docker-compose-services}{%
\subsection{Docker Compose Services}}

The \texttt{docker-compose.yml} defines all services needed for local
development:

\emph{(Mã nguồn chi tiết đã được lược bỏ để báo cáo ngắn gọn. Vui lòng
xem trong source code dự án)}

\hypertarget{starting-the-environment}{%
\subsection{Starting the Environment}}

\begin{Shaded}
\begin{Highlighting}[]
\CommentTok{\# Start all services in background}
\ExtensionTok{docker}\NormalTok{ compose up }\AttributeTok{{-}d}

\CommentTok{\# Check status}
\ExtensionTok{docker}\NormalTok{ compose ps}

\CommentTok{\# View logs}
\ExtensionTok{docker}\NormalTok{ compose logs }\AttributeTok{{-}f}\NormalTok{ postgres}
\ExtensionTok{docker}\NormalTok{ compose logs }\AttributeTok{{-}f}\NormalTok{ kafka}
\end{Highlighting}
\end{Shaded}

Expected output:

\begin{verbatim}
NAME                 STATUS
newsrag-postgres     Up (healthy)
newsrag-qdrant       Up
newsrag-zookeeper    Up
newsrag-kafka        Up (healthy)
newsrag-kafka-ui     Up
\end{verbatim}

\hypertarget{verifying-services}{%
\subsection{Verifying Services}}

\hypertarget{postgresql}{%
\subsubsection{PostgreSQL}}

\begin{Shaded}
\begin{Highlighting}[]
\CommentTok{\# Connect to database}
\ExtensionTok{psql} \AttributeTok{{-}h}\NormalTok{ localhost }\AttributeTok{{-}U}\NormalTok{ postgres }\AttributeTok{{-}d}\NormalTok{ newsrag}

\CommentTok{\# Or via Docker}
\ExtensionTok{docker}\NormalTok{ exec }\AttributeTok{{-}it}\NormalTok{ newsrag{-}postgres psql }\AttributeTok{{-}U}\NormalTok{ postgres }\AttributeTok{{-}d}\NormalTok{ newsrag}

\CommentTok{\# Verify tables created from warehouse.sql}
\ExtensionTok{\textbackslash{}dt}
\end{Highlighting}
\end{Shaded}

Expected tables:

\emph{(Mã nguồn chi tiết đã được lược bỏ để báo cáo ngắn gọn. Vui lòng
xem trong source code dự án)}

\hypertarget{qdrant}{%
\subsubsection{Qdrant}}

\begin{Shaded}
\begin{Highlighting}[]
\CommentTok{\# Check collections}
\ExtensionTok{curl}\NormalTok{ http://localhost:6333/collections}

\CommentTok{\# Expected: \{"result":\{"collections":[]\},"status":"ok","time":0.0...\}}
\end{Highlighting}
\end{Shaded}

\hypertarget{kafka}{%
\subsubsection{Kafka}}

\begin{Shaded}
\begin{Highlighting}[]
\CommentTok{\# List topics}
\ExtensionTok{docker}\NormalTok{ exec newsrag{-}kafka kafka{-}topics }\AttributeTok{{-}{-}bootstrap{-}server}\NormalTok{ localhost:9092 }\AttributeTok{{-}{-}list}

\CommentTok{\# Create news\_raw topic (if not auto{-}created)}
\ExtensionTok{docker}\NormalTok{ exec newsrag{-}kafka kafka{-}topics }\DataTypeTok{\textbackslash{}}
  \AttributeTok{{-}{-}bootstrap{-}server}\NormalTok{ localhost:9092 }\DataTypeTok{\textbackslash{}}
  \AttributeTok{{-}{-}create} \AttributeTok{{-}{-}topic}\NormalTok{ news\_raw }\DataTypeTok{\textbackslash{}}
  \AttributeTok{{-}{-}partitions}\NormalTok{ 3 }\AttributeTok{{-}{-}replication{-}factor}\NormalTok{ 1}
\end{Highlighting}
\end{Shaded}

\hypertarget{kafka-ui}{%
\subsubsection{Kafka UI}}

Open http://localhost:8080 in browser to view topics, messages, consumer
groups.

\hypertarget{python-environment-setup}{%
\subsection{Python Environment Setup}}

\emph{(Mã nguồn chi tiết đã được lược bỏ để báo cáo ngắn gọn. Vui lòng
xem trong source code dự án)}

\hypertarget{requirements.txt}{%
\subsubsection{requirements.txt}}

\emph{(Mã nguồn chi tiết đã được lược bỏ để báo cáo ngắn gọn. Vui lòng
xem trong source code dự án)}

\hypertarget{environment-configuration}{%
\subsection{Environment Configuration}}

\begin{Shaded}
\begin{Highlighting}[]
\CommentTok{\# Copy example and edit}
\FunctionTok{cp}\NormalTok{ .env.example .env}
\end{Highlighting}
\end{Shaded}

\hypertarget{env.example}{%
\subsubsection{.env.example}}

\emph{(Mã nguồn chi tiết đã được lược bỏ để báo cáo ngắn gọn. Vui lòng
xem trong source code dự án)}

\hypertarget{running-the-local-pipeline}{%
\subsection{Running the Local
Pipeline}}

\hypertarget{crawl-news-scrapy-rightarrow-kafka}{%
\subsubsection{\texorpdfstring{1. Crawl News (Scrapy \(\rightarrow\)
Kafka)}{1. Crawl News (Scrapy \textbackslash rightarrow Kafka)}}}

\begin{Shaded}
\begin{Highlighting}[]
\CommentTok{\# Terminal 1: Start Kafka consumer (writes to PostgreSQL)}
\ExtensionTok{python}\NormalTok{ consumer/consumer.py}

\CommentTok{\# Terminal 2: Run crawler}
\ExtensionTok{python}\NormalTok{ main.py }\AttributeTok{{-}{-}mode}\NormalTok{ crawl}
\end{Highlighting}
\end{Shaded}

Expected output:

\begin{verbatim}
2024-01-15 10:30:45 [INFO] Spider opened: news_sitemap
2024-01-15 10:30:46 [INFO] Crawling VnExpress sitemap: https://vnexpress.net/sitemap_news.xml
2024-01-15 10:30:47 [INFO] Found 150 URLs in sitemap
2024-01-15 10:30:48 [INFO] Parsing article: https://vnexpress.net/kinh-doanh/...
2024-01-15 10:30:49 [INFO] Sent to news_raw[0]@123
...
2024-01-15 10:45:30 [INFO] Closing spider (finished)
\end{verbatim}

\hypertarget{run-etl-postgresql-rightarrow-star-schema}{%
\subsubsection{\texorpdfstring{2. Run ETL (PostgreSQL \(\rightarrow\)
Star
Schema)}{2. Run ETL (PostgreSQL \textbackslash rightarrow Star Schema)}}}

\begin{Shaded}
\begin{Highlighting}[]
\ExtensionTok{python}\NormalTok{ main.py }\AttributeTok{{-}{-}mode}\NormalTok{ etl}
\end{Highlighting}
\end{Shaded}

Expected output:

\begin{verbatim}
2024-01-15 10:46:00 [INFO] ETL processing batch of 50 articles
2024-01-15 10:46:05 [INFO] Upserted dim_source: VnExpress (vnexpress.net)
2024-01-15 10:46:05 [INFO] Upserted dim_time: 2024-01-15 08:30:00+00:00
2024-01-15 10:46:05 [INFO] Upserted dim_content: GDP quý 4 tăng 6.5%...
2024-01-15 10:46:06 [INFO] Inserted fact_articles: 50 articles
2024-01-15 10:46:06 [INFO] Created 327 chunks
2024-01-15 10:46:06 [INFO] ETL batch complete: 50 articles processed
\end{verbatim}

\hypertarget{vectorize-star-schema-rightarrow-qdrant}{%
\subsubsection{\texorpdfstring{3. Vectorize (Star Schema \(\rightarrow\)
Qdrant)}{3. Vectorize (Star Schema \textbackslash rightarrow Qdrant)}}}

\begin{Shaded}
\begin{Highlighting}[]
\ExtensionTok{python}\NormalTok{ main.py }\AttributeTok{{-}{-}mode}\NormalTok{ vectorize}
\end{Highlighting}
\end{Shaded}

Expected output:

\begin{verbatim}
2024-01-15 10:46:30 [INFO] Loading embedding model: BAAI/bge-small-en-v1.5
2024-01-15 10:46:35 [INFO] Model loaded. Embedding dimension: 384
2024-01-15 10:46:35 [INFO] Creating Qdrant collection: news_chunks
2024-01-15 10:46:35 [INFO] Vectorized batch: 256 chunks
2024-01-15 10:46:38 [INFO] Vectorized batch: 71 chunks
2024-01-15 10:46:38 [INFO] Vectorization complete. Total chunks: 327
\end{verbatim}

\hypertarget{run-full-pipeline-automated}{%
\subsubsection{4. Run Full Pipeline
(Automated)}}

\begin{Shaded}
\begin{Highlighting}[]
\ExtensionTok{python}\NormalTok{ main.py }\AttributeTok{{-}{-}mode}\NormalTok{ full}
\end{Highlighting}
\end{Shaded}

This runs crawl \(\rightarrow\) etl \(\rightarrow\) vectorize in
sequence.

\hypertarget{test-rag-api-interactive}{%
\subsubsection{5. Test RAG API
(Interactive)}}

\begin{Shaded}
\begin{Highlighting}[]
\CommentTok{\# Start FastAPI server}
\ExtensionTok{uvicorn}\NormalTok{ search.api:app }\AttributeTok{{-}{-}reload} \AttributeTok{{-}{-}port}\NormalTok{ 8000}

\CommentTok{\# Test in browser}
\ExtensionTok{open}\NormalTok{ http://localhost:8000/docs}
\end{Highlighting}
\end{Shaded}

Or via CLI:

\begin{Shaded}
\begin{Highlighting}[]
\ExtensionTok{python} \AttributeTok{{-}c} \StringTok{"}
\StringTok{from search.engine import Pipeline}
\StringTok{p = Pipeline()}
\StringTok{result = p.ask(\textquotesingle{}Tóm tắt tin tức về kinh tế Việt Nam hôm nay\textquotesingle{})}
\StringTok{print(result.summary)}
\StringTok{print(f\textquotesingle{}Sources: \{result.total\}, Time: \{result.duration\_ms\}ms\textquotesingle{})}
\StringTok{"}
\end{Highlighting}
\end{Shaded}

\hypertarget{makefile-commands}{%
\subsection{Makefile Commands}}

\emph{(Mã nguồn chi tiết đã được lược bỏ để báo cáo ngắn gọn. Vui lòng
xem trong source code dự án)}

\hypertarget{troubleshooting}{%
\subsection{Troubleshooting}}

\begingroup
\small
\setlength{\tabcolsep}{3pt}
\renewcommand{\arraystretch}{1.15}
\begin{longtable}[]{@{}
  >{\raggedright\arraybackslash}p{(\columnwidth - 2\tabcolsep) * \real{0.4118}}
  >{\raggedright\arraybackslash}p{(\columnwidth - 2\tabcolsep) * \real{0.5882}}@{}}
\toprule\noalign{}
\begin{minipage}[b]{\linewidth}\raggedright
Issue
\end{minipage} & \begin{minipage}[b]{\linewidth}\raggedright
Solution
\end{minipage} \\
\midrule\noalign{}
\endhead
\bottomrule\noalign{}
\endlastfoot
\texttt{PostgreSQL\ connection\ refused} & Wait for healthcheck:
\texttt{docker\ compose\ ps} should show \texttt{(healthy)} \\
\texttt{Kafka\ not\ ready} & Wait 30-60s for Kafka to start; check
\texttt{docker\ compose\ logs\ kafka} \\
\texttt{Qdrant\ collection\ not\ found} & Run vectorize step first, or
create manually via API \\
\texttt{Model\ download\ fails} & Check internet; model caches at
\texttt{\textasciitilde{}/.cache/huggingface/} \\
\texttt{Port\ already\ in\ use} & Stop conflicting services; change
ports in docker-compose.yml \\
\texttt{CUDA\ out\ of\ memory} & Use CPU:
\texttt{export\ CUDA\_VISIBLE\_DEVICES=""} \\
\end{longtable}
\endgroup

\hypertarget{development-workflow}{%
\subsection{Development Workflow}}

\begin{verbatim}
1. Edit code (crawler/spiders/spider.py, etl/etl_warehouse.py, etc.)
2. Test locally: make crawl / make etl / make vectorize
3. Check logs: docker compose logs -f <service>
4. Verify in DB: psql -h localhost -U postgres -d newsrag
5. Commit changes
6. Build Docker image for AWS deployment (see 5.9-AWS-Deploy)
\end{verbatim}

\hypertarget{next-steps}{%
\subsection{Next Steps}}

After local development works: 1. \textbf{Build \& Push Docker Image}
--- See \href{5.9-AWS-Deploy/}{AWS Deployment Prep} 2. \textbf{Deploy to
AWS} --- \href{5.10-Fargate-Crawler/}{Fargate Crawler},
\href{5.11-Lambda-Consumer/}{Lambda Consumer},
\href{5.12-Lambda-ETL/}{Lambda ETL} 3. \textbf{Test RAG API} ---
\href{5.13-RAG-API/}{RAG API}

\begin{center}\rule{0.5\linewidth}{0.5pt}\end{center}

\textbf{Next:} \href{5.5-Crawler/}{Crawler Development}
